% 本文件由 LaTeX 排版 Agent 根据有效 translation.json 自动生成。
% author=Tertullian; work_id=0315; title=论基督的肉身

\chapter{论基督的肉身}

% structure: p0001 source=h1 target=omitted reason=page-title-already-rendered-by-parent

\EditorialNote{本书是作者为驳斥某些否认基督肉身真实性，或至少否认其与人类肉身同一性的异端者而写——他们担心，若承认基督肉身的真实性，便也必须承认他在肉身中的复活，因而也承认人死后肉身复活。}\par

% structure: p0003 source=h2 target=section reason=relative-heading-rank; base=h2

\section{第一章 本著作的总旨。异端者——\Person{马尔基昂}、\Person{阿佩莱斯}和\Person{瓦伦提诺}——想要反驳复活的教义，便通过否认基督的肉身而剥夺他经历任何这类变化的可能性。}

他们如此急切地想动摇在我们现代的撒杜塞人出现之前早已稳固建立的复活信仰，以至于甚至否认这盼望与肉身有任何关系，因而也有充分理由用疑难问题来困扰基督的肉身，仿佛它要么根本不存在，要么具有一种完全不同于人类肉身的性质。\Emph{因为他们不能不感到忧虑}：一旦确定\Emph{基督的肉身}是人类的，就会立即产生一个反对他们的推论，即那已经在基督内复活的肉身，无论如何也必须复活。因此，我们必须从同样的武器库中保卫我们的复活信仰，他们正是从那里获得毁灭的武器。让我们查考主的身体实质，因为关于他的神性，众人一致同意。有问题的是他的肉身。其真实性和性质是争论的焦点。它可曾存在？从何而来？它是什么样的？如果我们能证明它，我们就为自己的复活确立一条法则。\Person{马尔基昂}为了否认基督的肉身，也否认了他的诞生；或者，他否认他的肉身，以便否认他的诞生；因为，他当然害怕，他的诞生和肉身会相互证明对方的真实性，因为没有肉身就没有诞生，没有诞生就没有肉身。仿佛，在那种一切异端共通的自负驱使下，他也不可能很好地要么否认诞生而承认肉身——就像\Person{阿佩莱斯}（他先是\Person{马尔基昂}的门徒，后来成为叛离者）那样——要么既承认肉身又承认诞生，却以不同的方式解释它们，就像\Person{瓦伦提诺}那样，他在门徒和叛离\Person{马尔基昂}方面与\Person{阿佩莱斯}相似。无论如何，那个认为基督的肉身是幻象的人同样能够把他的诞生当作幻影；因此，童贞女的受孕、怀孕、生育，以及后来她婴儿期的整个过程，都必须被视为拟似的。\Emph{这些属于基督诞生的事实}，将逃过那些未能全面把握其肉身之观念的眼睛和感官。\par

% structure: p0005 source=h2 target=section reason=relative-heading-rank; base=h2

\section{第二章 \Person{马尔基昂}妄想抹去基督诞生记录，被斥为如此骇人听闻的异端。}

\Person{加俾额尔}明确地宣告了诞生。\BibleRef{路 1:26} 但他与造物主的天使有何关系？童贞女腹中的怀孕也明确地摆在我们面前。但他与造物主的先知\Person{依撒意亚}有何关系？他将不耐烦地等待，因为他突然（没有任何预言宣告）将基督从天上带下来。\Quote{滚开，}他说，\Quote{那个永无止境的邪恶的\Person{凯撒}课税，以及简陋的旅馆，肮脏的襁褓，和坚硬的马厩。\BibleRef{路 2:1} 我们一点也不在乎那队在主夜间赞美主的众天军。\BibleRef{路 2:13} 让牧人更好地照看他们的羊群，\BibleRef{路 2:8} 让智者省下他们长途跋涉的双腿；\BibleRef{玛 2:1} 让他们自己留着黄金。\BibleRef{玛 2:11} 也让\Person{黑落德}改改他的作风，这样\Person{耶肋米亚}就不会因他而夸口。也饶了婴孩不受割礼，好让他免去疼痛；也不要将他带进圣殿，免得加重父母奉献的花费；\BibleRef{路 2:22} 也不要将他交给\Person{西默盎}，免得老人在临终时伤心。\BibleRef{路 2:25} 也让那个老妇人闭上嘴，免得她给婴儿施了妖术。\BibleRef{路 2:36}} 我想，你就是这样，\Person{马尔基昂}，胆敢抹去基督历史的原始记录，好让他的肉身失去其真实性的证据。但是，请问，你凭什么理由这样做？给我看看你的权威。如果你是一个先知，就预言一件事；如果你是一个宗徒，就在公开场合发表你的信息；如果是跟随宗徒的人，就与宗徒同心；如果你只是一个（普通）基督徒，就相信那传给我们的；如果你什么都不是，那么（我有最好的理由说）你就别再活了。因为实际上你已经死了，既然你不是基督徒，因为你不相信那使人成为基督徒的信仰——不，你死得更厉害，你越不是基督徒；你曾经是基督徒后又堕落了，因为你拒绝了从前相信的，正如你自己在某一封信中承认的，你的追随者也不否认，而我们的弟兄可以证明。因此，拒绝了你曾经相信的，你就完成了拒绝的行为，因为现在不再相信了：但是，你停止相信这个事实并没有使你对信仰的拒绝变得正当；相反，通过你的拒绝行为，你证明了你在这个行为之前所相信的是另一种性质。你所相信的是另一种性质，那传下来的正是你所相信的。另一方面，那传下来的是真的，因为是由那些有责任传递它的人传递的。所以，当你拒绝那传下来的时，你拒绝了那真的。你这样做没有权威。然而，我们已经在一篇别的论文中更充分地利用了这些反对一切异端的规则。我们在那篇长篇大论之后重复它们是没有必要的，因为我们要问的是你为什么形成基督没有诞生的见解。\par

% structure: p0007 source=h2 target=section reason=relative-heading-rank; base=h2

\section{第三章 基督的诞生既是可能的又是相称的。关于基督外表肉身的异端见解具有欺骗性且有损天主尊严，甚至依照\Person{玛尔基翁}的原则亦是如此}

既然你认为这取决于你自己任意的选择，你就必定认为诞生对天主来说既是不可能的，或是对祂不相称的。然而，对天主而言，除了祂不愿意的，没有什么是做不到的。那么，让我们思考：祂是否愿意被生（因为如果祂愿意，祂也有能力，并且祂被生了）。我简洁地提出论证。如果天主不愿意被生，无论原因如何，祂就不会以人的样式显现自己。如今，谁看到一个人，会否认他被生了呢？因此，\Emph{天主}不愿意成为的，祂绝不会愿意貌似是那样。当一件事令人反感时，就连它的概念也会被拒绝；因为一件事存在或不存在并无区别，如果，当它不存在时，它却被假定存在。当然，最重要的是，不应有任何虚假（\Emph{或假装}）归于实际上不存在的东西。但你说，祂自己的意识（对祂本性的真相）对祂来说足够了。如果有人因为看见祂是一个人而认为祂被生了，那是他们的事。然而，假设祂确实被生了，祂会以更大的尊严和一致性来维持人的性情；\Emph{因为}如果祂没有被生，祂就不能承担那性情，而不损害祂的那种意识，你们方面将那种意识归因于\Emph{祂}的自信，即尽管没有被生，仍能维持被生的性情，甚至违背祂自己的意识！为什么（我想知道）基督在完全清楚祂实际上是怎样的时，要表现出祂不是的那样，竟如此重要？你不能表达任何担忧：如果祂被生并真正披上人的本性，祂就会不再是天主，在成为祂所不是的同时失去祂所是的。因为天主不会有失去祂自己的状态和状况的危险。但你说，我否认天主真的变成了人，以致被生并赋予血肉之躯，理由是：一个没有终结的存在也必然是不变的。因为变成别的东西结束了先前的状态。因此，对一个不会终结的存在来说，变化是不可能的。无疑，会变化的事物的本性受这条规律支配：它们在其发生变化的状况中没有恒常性，并且由于缺少恒常性而终结，在变化过程中失去它们先前所是。但没有什么能与天主相等；祂的本性与万物的状况不同。那么，如果那些与天主不同、天主也与它们不同的事物，在发生变化时失去它们原有的存在，那么天主与万物之间的区别将在于什么，除非在于祂拥有与它们相反的能力——换句话说，天主可以变成一切状况，却仍然保持祂的本然？在任何其他假设下，祂就会与那些变化后失去原有存在的东西同一层次；当然，祂在任何其他方面都不与它们同等，正如祂当然不在它们本性的变化结果上同等。你有时读到并相信造物主的天使变成了人的形状，甚至还带着如此真实的躯体，以至亚巴郎曾洗他们的脚，\EditorialNote{创世纪第十八章}，罗特被他们从索多玛人手中救出；\EditorialNote{创世纪第十九章}；而且，一位天使用身体与一个人如此激烈地摔跤，以致那人想要被放开，因为他被紧紧抓住。\EditorialNote{创世纪第三十二章}那么，是否允许低于天主的天使们，在变成人的形体之后，却仍然保持为天使？而你要剥夺天主——他们的上主——的这种能力，仿佛基督在真正取人性之后不能继续是天主？或者，那些天使只是以血肉的幻影出现？然而你没有勇气这样说；因为如果在你的信仰中，造物主的天使与基督处于同样的状况，那么基督就属于与那些天使相同的那位天主，那些天使在状况上类似基督。如果你没有在一些场合故意拒绝，在另一些场合篡改那些与你意见相悖的圣经，你就早被\BibleBook{若望}{福音}在这件事上驳倒了，因为它宣告圣神降在鸽子形体中，并落在主身上。\BibleRef{玛 3:16}当那圣神处于这种状况时，祂是真的鸽子，正如祂也是神；祂并没有因取用外来的实体而毁坏祂自己的实体。但你会问，在圣神返回天上之后，鸽子的身体变成了什么，天使的情况也类似。他们的离去与他们显现的方式相同。如果你看到他们从无中产生是如何实现的，你也会知道他们回归于无的过程。如果最初的一步是不可见的，最后的一步也是如此。然而，他们的身体实体仍有坚固性，无论那身体变得可见的力量是什么。所记载的不能不如此。\par

% structure: p0009 source=h2 target=section reason=relative-heading-rank; base=h2

\section{第四章 天主在其圣子降生成人中的尊荣得以辩护。\Person{玛尔基翁}对人的肉身的诋毁既矛盾又不敬。基督已洁净了肉身。天主的愚拙是最有智慧的}

所以，既然你不认为取得一个身体是不可能的，或对天主的本性有害，那么你只好把它当作与祂不相称而加以拒绝和指责。来吧，从头开始，从诞生本身开始，谴责子宫内生殖元素的不洁，液体和血液的污秽凝结，以及肉体从那泥潭中生长九个月之久。描述那日益增大的子宫，沉重、麻烦、甚至睡中也不安宁，在厌恶和欲望的感觉中变化不定。现在同样抨击产妇的羞耻本身——然而，鉴于那危险，这理应受到尊敬，或者出于（奥秘的）本性被视为神圣。当然，你也对那带着子宫的窘迫而出生到世上的婴孩感到恐惧；当然，你同样在婴孩洗浴之后也厌恶他，当他被裹在襁褓中，被反复涂油美化，被保姆的宠爱逗笑时。\Person{马尔西翁}啊，你唾弃这庄严的自然过程；然而，你是怎样出生的呢？你憎恨出生时的人；那么，你又以何种方式爱任何人呢？当然，你离开教会和基督的信仰时，你并不爱你自己。但没关系，如果你与自己不睦，或者你的出生方式与别人不同。无论如何，基督甚至爱了那个在母腹中凝聚于所有不洁之中的人，甚至爱了那个从那子宫中出生的人，甚至爱了那个在保姆的微笑中受哺育的人。为了他，祂降下（从天），为了他，祂宣讲，为了他，祂\Quote{谦抑自己，听命至死，且死在十字架上。}\BibleRef{斐 2:8} 当然，祂爱了那个祂以如此大代价救赎的人。如果基督是造物主的\Emph{儿子}，祂爱自己的（受造物）是正义的；如果祂来自另一位神，祂的爱就过分了，因为祂救赎了一个属于别人的受造物。那么，爱人的祂也爱人的诞生，也爱人的肉体。除非通过那使得存在物得以存在的条件，没有东西能被爱。要么除去诞生，然后向我们展示\Emph{你的}人；要么撤去肉体，然后向我们呈现天主所救赎的那个存在——因为正是这些条件构成了天主所救赎的人。而\Emph{你}，却要将这些条件变成那被救赎者脸红的原因，并认为它们与祂不相称，而祂若不爱它们，就不会救赎它们？祂以第二次从天上的出生，将我们的出生从死亡中改造；祂将我们的肉体从各种折磨的疾病中恢复；当它患麻风病时，祂洁净其污点；当它失明时，祂重新点燃其光明；当它瘫痪时，祂更新其力量；当它附魔时，祂驱魔；当它死亡时，祂重新赋予生命——那时\Emph{我们}还会羞于承认它吗？当然，如果祂选择只是从一个动物出生，并披戴一只野兽（野生或驯服）的身体来宣讲天国，我想你的指责会立刻用这个异议与祂相遇：\Quote{这对天主是可耻的，这与天主子不相称，简直是愚蠢。} 没有其他原因，只是因为有人如此判断。这\Emph{是}当然愚蠢的，如果我们按自己的概念判断天主。但是，\Person{马尔西翁}，请仔细考虑这经文，如果你没有删掉它：\Quote{天主拣选了世上愚拙的，叫那有智慧的羞愧。}\BibleRef{格前 1:27} 那么，这些愚拙的是什么呢？是使人皈依真天主的敬拜，拒绝错误，全部公义、贞洁、仁慈、忍耐和纯洁的训练？这些东西当然不是\Quote{愚拙的。} 那么再查考一下祂说的是什么，当你以为已经发现了它们时，你会发现有什么像相信一位诞生的天主、由童贞女所生、且是肉身的本性、祂沉浸于前述本性的所有羞辱中那样\Quote{愚拙}吗？但有人会说：\Quote{这些不是愚拙的事；必是别的天主所拣选以扰乱世界智慧的事。} 然而，按照世界的智慧，相信\Person{朱庇特}变成一头公牛或一只天鹅（如果我们听信\Person{马尔西翁}）比相信基督真的成为一个人更容易。\par

% structure: p0011 source=h2 target=section reason=relative-heading-rank; base=h2

\section{第五章 基督真实地在肉身中生活和死亡。祂在地上人类生活的种种事件，以及对\Person{马尔西翁}幻象论式模仿的驳斥}

当然，还有其他同样愚昧（如基督的诞生）的事，涉及天主的屈辱和苦难。或者，让他们称被钉死的神为\Quote{智慧}。但马尔基昂也会对这条\OriginalTerm{道理}{doctrine}动刀，甚至更有理由。因为哪一样更不适合天主，哪一样更易让人脸红羞愧：是天主应当诞生，还是祂应当死去？是祂应当承载肉身，还是十字架？是受割损，还是被钉？是睡在摇篮，还是躺在棺木？是放在马槽，还是坟墓？\Emph{说什么}\Quote{\Emph{智慧}！}你若拒绝相信这些，你倒会显得更有\Quote{智慧}。但毕竟，除非你成为世界的\Quote{愚人}，相信\Quote{天主的愚昧事}，你就不会\Quote{有智慧}。那么，你从基督身上割去了所有苦难，认为祂只是个幻象，无法经历它们？我们上面说过，祂可能经历了想象中出生和童年的虚幻嘲弄。但你，这个谋杀真理的人，立刻回答我：天主不是真的被钉死了吗？而且真的被钉死后，不是真的死了吗？并且真的死后，不是真的复活了吗？保禄\Quote{决定在我们中间只知道耶稣，并被钉在十字架上的祂}，岂不是假的？\BibleRef{格前 2:2}祂被埋葬，岂不是假教导的？祂复活，岂不是假灌输的？因此，我们的信仰也是假的。我们一切从基督所希望的，都将成为幻象。哦，你这些最无耻的人，竟宣告杀害天主的凶手无罪！因为如果基督真的什么也没受，那祂从他们那里什么也没受。请放过世界唯一的希望吧，你正在摧毁我们信仰必不可少的羞辱。凡是不适合天主的，对我都是增益。倘若我不以我的主为耻，我就安全。祂说：\Quote{凡以我为耻的，我也将以他为耻。}我看不到别的可耻的事，能使我对可耻的事有好的意义上的无耻，并因我对可耻的轻视而成为幸福的愚昧。天主子被钉死了；我不以为耻，因为人必须为它而羞耻。天主子死了：这是完全值得相信的，因为它是荒谬的。祂被埋葬了，又复活了：这事实是确定的，因为它是不能的。但这一切在祂身上如何真实，如果祂本身不是真实的——如果祂确实没有那可被钉死、可死、可埋、可复活的东西呢？\Emph{我是说}这肉身充满血液，由骨头建造，由神经交织，由血管缠绕，\Emph{这肉身}知道如何出生，如何死亡，无疑是人的，由人而生。因此，它在基督内必是必朽的，因为基督是人，也是人子。否则，基督为什么是人，又是人子，如果祂没有人的任何东西，也不从人来的任何东西？除非要么人与肉身不同，要么人的肉身从人以外的源头而来，要么玛利亚不是人，要么马尔基昂的人如同马尔基昂的神。否则，基督不能被称为没有肉身的人，也不能是没有人类父母的人子；正如祂没有天主的圣神就不是天主，没有天主为父就不是天主子。这样，两种本质的性质将祂显示为人与天主——一方面被生，另一方面非生；一方面属血肉，另一方面属神；一方面软弱，另一方面极其强大；一方面会死，另一方面永生。这两种状态——神圣与人性——的性质，同样真实地断言于两种本性，既相信关于圣神也关于肉身。圣神的德能证明祂是天主，祂的苦难证实了人的肉身。如果祂的德能不是没有圣神，同样，祂的苦难也不是没有肉身。如果祂的肉身及其苦难是虚构的，那么圣神及其一切德能也是假的。因此，你为什么用谎言分割基督？祂完全是真理。相信我，祂宁愿被生，也不愿在任何部分假装——甚至确实对自己损害——祂带着一个没有骨骼而硬化、没有肌肉而坚实、没有血液而血淋、没有皮肤外衣而穿衣、没有食欲而饥饿、没有牙齿而吃东西、没有舌头而说话，使祂的话语通过想象的声音对耳朵成为幻象。当然，复活后也是幻象，当祂伸出手和脚让门徒查看时，祂说：\Quote{你们看看我的手和脚，就知道是我本人；因为鬼魂没有肉身和骨头，你们看我是有的。}\BibleRef{路 24:39}毫无疑问，手、脚和骨头不是鬼魂所拥有的，而\Emph{只有}肉身。你怎么解释这话，马尔基昂，你告诉我们耶稣只来自最卓越的天主，祂既单纯又良善？你看祂\Emph{反而}欺骗、迷惑、戏弄所有人的眼睛和所有感官，以及与祂的接触和触摸！你更该把基督带下来，不是从天上，而是从某个戏班子，不是作为天主而又是人，而只是作为人，一个魔术师；不是作为我们救恩的大司祭，而是作为表演中的变戏法者；不是作为使死人复活者，而是作为误导活人者——除非，如果祂是魔术师，祂必然有出生！\par

% structure: p0013 source=h2 target=section reason=relative-heading-rank; base=h2

\section{第六章 阿佩莱斯的道理被驳斥：基督的身体并非由星质构成，并非受生。出生与死亡是相互关联的情形；在基督的情况下，祂的死证实了祂的生}

然而，本都（\Place{本都}）异端者的一些门徒，被迫比他们的老师更为明智，便承认基督有真实的肉身，但这并不影响他们否认祂的诞生。他们说，祂可能拥有一个从未出生的肉身。于是，我们正如俗语所说，\Quote{才出油锅}，\Quote{又入火坑}——从马尔强（\Person{马尔强}）到阿贝肋（\Person{阿贝肋}）。此人首先在肉体上从马尔强的原则堕落，与一个女人纠缠，随后在精神上因童贞斐禄默纳（\Person{斐禄默纳}）而沉船，从\Emph{那时}起便开始宣讲基督的身体是坚实的肉身，但并未出生。对于这个斐禄默纳的天使，宗徒将以他当时预言时同样的语气回答说：\BibleQuote{即使从天上来的天使，给你们宣讲与我们先前所宣讲的福音不同的福音，当受诅咒。}\BibleRef{迦 1:8}现在，我们必须抵制刚才所指出的论据。他们承认基督确实有一个身体。这身体来自何种材料，如果不是来自祂显现时所借以显现的那种同类东西，又是来自何处？如果祂的身体不是肉身，身体从何而来？如果肉身不是出生的，肉身从何而来？因为，要成为肉身，必须经历这种诞生。他们声称，祂从星辰和上界物质借来了肉身。他们断言一条确定的原则：一个没有出生的身体不足为奇，因为天使也曾在没有子宫干预的情况下，以肉身显现于我们当中。当然，我们承认已有此类事实的记载。但是，持不同规则的信德，如何能够从它所攻击的信德中借用材料来论证自己的论据呢？它与梅瑟（\Person{梅瑟}）有什么关系？它已经拒绝了梅瑟的天主。既然天主不同，属于祂的一切也必然不同。然而，让异端者总是使用那位天主的圣经吧，他们同样享用着祂的世界。这事实将反作用于他们，作为见证来审判他们，因为他们从\Emph{祂}的范例中维持自己的亵渎之言。然而，真理无需提出此类异议就能轻易得胜。因此，当他们以天使为范例提出基督的肉身，宣称它虽未出生却仍是肉身时，我希望他们比较基督与天使两者以肉身显现的原因。没有天使曾降临为要钉在十字架上，尝死味\Emph{并}复活。既然天使从未有成为肉身这样的理由，你就有原因知道他们为何不经出生就取了肉身。他们并非来受死，因此他们也不为出生。然而，基督被派遣来受死，就必然也要出生，好使祂能死；因为只有出生的才会死。出生与死亡之间有着相互的对照。使我们死亡的法律正是我们出生的原因。既然基督是因那承受死亡的条件而死，而那承受死亡的也是出生的，那么结果——不，这是必然的前提——祂也必须因那承受出生的条件而出生；因为祂必须服从那从出生开始、以死亡结束的同一条件而死。祂不应以应死为借口而不出生。然而主自己在那时就曾在天使中向亚巴郎（\Person{亚巴郎}）显现，并未出生，却无疑具有肉身，这是因为上述原因的差异。但是你不能承认这一点，因为你没有接受那位基督，祂当时就在那尚未出生的肉身形态中预演了如何与人交谈、解放人类并审判人类，因为祂那时还不打算死，直到双方（出生和死亡）都（由预言）先被宣布。那么，让他们向我们证明那些天使从星辰取得肉身。如果他们无法证明，因为圣经没有记载，那么基督的肉身也不会由此而来，尽管他们借用了天使的范例。显然，天使所拥有的肉身并非本性的；他们的本性是精神实体，尽管在某种意义上是可形的，但他们可以变形为人的形象，并暂时出现与人交往。既然我们没有被告诉他们从何处获得肉身，我们就应心里不疑惑：天使之力的特性是能够从无物质材料中给自己取形体。你说，何况从某些物质材料中取形体呢？这很对。但没有证据，因为圣经没有说。再者，那些能将自己形变成本来不是之物的人，怎能没有能力从无物中如此行呢？如果他们成为他们本不是之物，为何不能从非存在之物\Emph{如此}成为呢？但那不存在而存在之物，是从无中\Emph{造成}的。这就是为何无需追问或演示他们身体后来如何。从无中来的，归于无。能够把\Emph{自己}变成肉身的，也有能力把\Emph{无本身}变成肉身。改变本性比创造物质更伟大。但即使必须\Emph{假设}天使是从某种物质材料取得肉身，也更有理由相信是从地上的物质，而非某种天上的物质，因为他们的肉身显然具有可触的地上性质，以地上的食物为生。假设即使天上的\Emph{肉身}也以地上的食物为生，虽然它本身不是地上的（正如地上的肉身确实以天上的食物为生，虽然它没有天上的性质——我们读到玛纳是民众的食物：\Emph{圣咏作者}说：\BibleQuote{人吃了天使的食粮}），但这并不因此损害主肉身的独特状况，因为祂有不同的目的。因为那要真正成为人、甚至至于死的一位，必须穿戴上那属于死亡的肉身。而那属于死亡的肉身，是出生所先行的。\par

% structure: p0015 source=h2 target=section reason=relative-heading-rank; base=h2

\section{第七章 解释主关于祂的母亲和兄弟的提问。答复\Person{阿佩莱斯}和\Person{马尔基翁}的吹毛求疵的异议，他们藉此支持他们否认基督诞生的真实性。}

但每当关于诞生的争论兴起时，所有拒绝它的人——因为它造成了有利于基督肉身实在性的推定——都故意否认天主自己曾降生，其理由是祂曾问：\BibleQuote{谁是我的母亲？谁是我的兄弟？}所以，让\Person{Apelles}听听我们在那篇小作品中给\Person{Marcion}的答复；在那作品中，我们曾向祂自己（最喜爱的）福音挑战以求证明，即那番话（主的）的具体情境应被考虑。首先，没有人会告诉祂祂的母亲和兄弟们站在外面，除非他确定祂有母亲和兄弟，而且他们正是他当时所宣布的那些人——这些人他以前就认识，或者当时当场发现了；尽管异端者已经从福音中删除了这段经文，因为那些钦佩祂教义的人说，他们非常熟悉祂所谓的父亲若瑟木匠、祂的母亲玛利亚、祂的兄弟和祂的姊妹。但他们是怀着试探祂的意图，才向祂提起祂所没有的母亲和兄弟。圣经对此只字未提，尽管在其他场合，当有人以试探的方式对祂行事时，圣经并非沉默。\BibleQuote{看，一个法学士站起来，试探祂。} \BibleRef{路 10:25} 又在另一处：\BibleQuote{法利塞人也来到祂那里，试探祂。} 谁能阻止此处也表明这是出于试探祂的意图呢？我不承认你们离开圣经自己所提出的那些。那么，应当为试探提出某种时机。他们认为在祂内有什么需要试探的呢？当然，是祂是否曾降生的问题？因为如果祂的回答否认了这一点，那么当宣布试探时，这一点就可能暴露出来。然而，当试探旨在发现那个因可疑而引发试探的要点时，没有试探会如此突然地降临，以至于没有先提出一个问题，而这个问题在引起怀疑的同时迫使试探发生。如今，既然基督的降生从未成为疑问，你们怎能争辩说他们藉着试探意图查明一个他们从未怀疑过的问题呢？此外，如果祂必须在诞生这方面受试探，这当然不是合适的方式——即宣布那些即使假定祂已诞生，也可能并不存在的人。我们都已降生，但并非人人都有兄弟或母亲。祂更可能有父亲而非母亲，更可能有叔伯而非兄弟。因此，关于祂诞生的试探是不合适的，因为完全可以不提祂的母亲或兄弟而设计试探。更可信的是，他们确信祂既有母亲也有兄弟，就试探祂的天主性而非祂的降生，即祂在里面是否知道外面的事；他们藉着虚假宣布那些并不在场的人在场来考验祂。但试探的计谋可能这样被挫败：或许祂知道那些被他们宣布为\Quote{站在外面}的人，实际上因生病、事务或祂当时所知的旅行而不在场。没有人会用一种他知道自己可能承受试探之耻的方式去试探（别人）。因此，既然没有合适的试探时机，宣布祂的母亲和兄弟真的出现了，就恢复了其自然性。但有理由认为，\Emph{基督的}回答暂时否认了祂的母亲和兄弟，正如连\Person{Apelles}也可能得知的那样。\BibleQuote{主的兄弟尚未相信祂。} \BibleRef{若 7:5} 这记载在马西昂时代之前就已出版的福音中；同时缺乏祂的母亲坚持与祂同在的证据，尽管玛尔大和其他玛利亚们经常服侍祂。的确，就在这段经文中，他们的不信是明显的。耶稣正在教导生命之道，宣讲天主的国，并\Emph{且}积极致力于医治身体和灵魂的软弱；但与此同时，当陌生人都专注于祂时，祂最近的亲属却缺席。不久他们出现了，却站在外面；他们不进去，因为（的确）他们对里面所进行的事不重视；他们甚至不等候，似乎他们能贡献比祂正热切做的事更必要的东西；但他们宁愿打断祂，希望把祂从伟大的工作中叫走。现在，我问你，\Person{Apelles}，或者你\Person{Marcion}，请问你们：如果你们碰巧在看戏剧，或对赛跑或赛车下了赌注，却被这样的消息叫走，你们难道不会喊出：\Quote{母亲和兄弟对我算什么？} 难道基督在宣讲并彰显天主、满全法律和先知、并\Emph{且}驱散漫长前代黑暗的同时，不正是恰当地使用了同样的措辞，以打击那些站在外面之人的不信，或摆脱那些想叫祂离开工作之人的纠缠吗？然而，如果祂本意是否认自己的降生，祂本会找到非常不同的地方、时间和方式来表达自己，而不会用那些既有母亲又有兄弟的人也可能说出的话。当义愤地否认自己的父母时，一个人并非否认\Emph{他们的存在}，而是斥责\Emph{他们的过错}。此外，祂给了别人优先权；既然祂显示他们得到这一恩宠的理由——正是因为他们听了（天主的）话——祂就指明了在何种意义上祂否认了祂的母亲和兄弟。因为无论祂在何种意义上将那些依附祂的人接纳为自己的，祂都在同一意义上将那些疏远祂的人否认为自己的。基督也习惯竭尽所能去行祂所吩咐别人的。那么，如果祂在教导别人不要将母亲、父亲或兄弟看得如同天主的话那样重要时，自己却在母亲和兄弟被宣告时就离开了天主的话，这会是多么奇怪啊！因此，祂否认了自己的父母，是在祂教导我们为了天主的工作而否认自己父母的意义上。但还有一种观点：在被弃绝的母亲身上，有着会堂的预象，而在不信的兄弟身上，有着犹太人的预象。在他们身上，以色列留在外面，而紧靠基督在内的新门徒，听到并相信，代表了教会，祂称教会为更优越意义上的母亲和更值得的兄弟关系，同时摒弃了血亲关系。确实，正是出于同一意义，祂也回应了那（某位妇女的）呼喊，没有否认祂母亲的\Quote{子宫和乳房}，而是指出那些\Quote{听天主的话而遵行的人}更有福。\par

% structure: p0017 source=h2 target=section reason=relative-heading-rank; base=h2

\section{第八章 \Person{阿佩莱斯}与其追随者不悦我们地上的身体，便赋予基督一种更纯净的身体；基督如何甚至在祂地上的肉身中也是属天的}

仅凭这些经文——\Person{阿佩莱斯}与\Person{马西昂}似乎在按整个未受玷污之福音的真理加以解释时，将其主要依靠放在了上面——本已足以通过捍卫基督的诞生来证明祂属人的肉身。但由于\Person{阿佩莱斯}那宝贵的群体过分强调肉身的可耻状况，他们认为肉身被那火性的邪恶作者所败坏之灵魂充满，因此不配于基督；并且因为他们因此认为一种星辰性的实体适合祂，我必须在他们自己的立场上反驳他们。他们提到某位大名鼎鼎的天使创造了我们这个世界，并在创造之后为他的工作后悔。这点我们实际上已在单独的一章中讨论过；因为我们写了一篇小作品反对他们，\Emph{论及}那拥有基督的、为如此行动所需之精神、意志和权能者，是否可能做出任何需要悔改之事的问题，因为他们以\Quote{亡羊}的比喻来描述那\Emph{所述的}天使。因此，按照其创造者悔改的证据，世界必定是错误的；因为一切悔改都是对过错的承认，而且它若非因过错根本不存在。现在，如果世界是过错，如同身体一样，那么其各部分也必是如此——同样是有过错的；同样，天上及其天体（内容），以及一切从中受孕并产生的事物，也必如此。而\BibleQuote{坏树必结坏果子。}\BibleRef{玛 7:17}因此，基督的肉身若由天体元素构成，便是由有瑕疵的材料组成，因其有罪的起源而有罪；所以它必定是那实体的一部分——那实体他们因有罪而不屑于用来装扮基督——换言之，就是我们自己的实体。那么，既然在可耻方面没有差别，就让他们要么为基督设计一种更纯净的实体，既然他们不喜悦我们的实体；要么也承认这个实体，即使天上的实体也不会比这更好。我们读到的原话是：\BibleQuote{第一个人是出于地，属于土；第二个人是主，自天而来。}\BibleRef{格前 15:47}然而，这段经文与实体的任何差异无关；它只是将第一个人亚当那曾经\Quote{属土}的肉身实体，与第二个人基督那\Quote{属天}的圣神实体相对照。而且，这段经文完全将属天之人归于圣神而非肉身，以至于那些与之相比的人显然成为属天的——当然是因着圣神——甚至在这\Quote{属土的肉身}中。现在，既然基督甚至在肉身方面也是属天的，那些在肉身方面不是属天的人就不能与祂相比。因此，如果他们成为属天的，如同基督也是属天的一样，却带着\Quote{属土的}肉身实体，那么这一事实所肯定的结论是，基督自己也是属天的，但是在\Quote{属土的}肉身中，如同那些与祂同等的人一样。\par

% structure: p0019 source=h2 target=section reason=relative-heading-rank; base=h2

\section{第九章 基督的肉身完全自然，如同我们自己的肉身；仔细查看之下，异端者归给祂的超自然特征一个也没有发现}

我们迄今所遵循的原则是，从某物衍生而来的任何东西，无论它与所从出的东西有多么不同，都不至于无法暗示其来源。没有一种物质质料不带有其自身起源的见证，无论它经历了多么巨大的性质变化。就以我们这身体为例，它由地上的尘土形成，这是一个连外邦人的寓言中也流传的真理；它确凿地见证了自己源于土和水两种元素——前者通过其肉，后者通过其血。虽说在品质的外观上存在差异（换句话说，从某物而来的东西在发展中是不同的），但归根结底，血不就是红色的液体吗？肉不就是一种特殊形态的土吗？思考各自的品质：肌肉如同土块，骨骼如同石头，乳腺如同卵石。看神经的紧密连接如同根的繁殖，血管的分支蜿蜒如同溪流，覆盖我们的茸毛如同苔藓，头发如同草，骨髓中的珍藏如同肉的矿藏。所有这些属地的起源的标志都在基督内；正是它们掩盖了祂作为天主子的事实，因为祂被看做是人，别无其他原因，只因为祂存在于人的身体实体中。否则，请给我们指出祂身上某种从大熊、昴星、毕星窃取来的天上的物质吧。那么，我们列出的这些\Emph{特征}就是大量的证据，证明祂的肉是\OriginalTerm{属地的肉}{earthy flesh}，如同我们的肉一样；但我发现不了任何新奇或奇怪的东西。实际上，人们正是从祂的言语和行动、祂的教导和奇迹，尽管感到惊奇，才承认基督是人。但如果祂身上有一种奇迹般获得的（从星辰来的）新奇的肉，那当然是众所周知的。然而，实际情况是，正是祂属地的肉的普通状况使得祂的一切其他方面显得神奇，正如他们所说：\Quote{这人从哪里得了这智慧和这些异能？}\BibleRef{玛 13:54}连那些轻视祂外貌的人也是这样说的。祂的身体甚至没有达到人的美丽，更谈不天上的荣耀。即使先知从未给我们任何关于祂卑微外貌的信息，祂所受的苦难和所忍受的凌辱本身也足以表明这一切。苦难证实了祂的人的肉，凌辱证明了其卑贱的状况。如果\Emph{基督}的身体具有不寻常的本性，谁会敢用指尖触碰它？或者向祂脸上吐唾沫，如果它没有（因其卑贱）招致呢？为何谈论属天的肉，当你没有提供任何理由来支持你的关于天的理论？为何否认它是属地的，当你完全有理由知道它是属地的？祂在魔鬼的\Emph{试探}下饥饿；祂在撒玛黎雅妇人那里口渴；祂为拉匝禄哭泣；祂在死亡面前颤抖（因为\Quote{肉是软弱的}\BibleRef{玛 26:41}）；最后，祂流出了自己的血。这些，我想，是天上的标志吗？但是，我请问，如果祂那肉中闪耀着任何天上的优越，祂怎能照我所描述的那样遭受轻蔑和苦难呢？因此，我们由此得到一个令人信服的证据：在祂的肉中没有任何属天的事物，因为它必须能够受轻蔑和苦难。\par

% structure: p0021 source=h2 target=section reason=relative-heading-rank; base=h2

\section{第十章 驳斥另一类异端者。他们声称基督的肉是更精细的质料，是\OriginalTerm{灵魂的}{Animalis}，由灵魂构成。}

我现在转向另一类异端者，他们在自己的想象中同样聪明。他们断言基督的肉由灵魂构成，祂的灵魂成了肉，以致祂的肉是灵魂；并且正如祂的肉出于灵魂，同样祂的灵魂也出于肉。但在这里，我又必须有一些理由。如果为了拯救灵魂，基督在自己内取了一个灵魂，因为灵魂除非祂拥有它，就不能得救，我看不出祂在穿衣时，为什么要把那肉造成灵魂的，好像祂除了使灵魂变成肉之外，就无法拯救灵魂似的。因为当祂拯救\Emph{我们}这些不仅不是肉、甚至与肉不同的灵魂时，祂更能确保祂自己所取的那个灵魂（即使它也不是肉）的救恩。此外，既然他们断言一个主要信条：基督来不是为了拯救肉，而只是为了我们的灵魂，那么首先，多么荒谬的是：祂意图只拯救灵魂，却把灵魂造成那种祂无意拯救的身体物质！其次，如果祂承担了藉着祂所携带的东西来拯救我们的灵魂，那么祂在那灵魂中所携带的，应当携带我们的\Emph{灵魂}，即一个与我们状况相同的灵魂；而我们的灵魂在其隐秘的本性中的任何状况，当然不是肉。然而，祂所拯救的不是我们的灵魂，如果祂自己的灵魂是肉的话；因为我们的灵魂不是肉。现在，如果祂没有拯救我们的灵魂，理由是祂所拯救的是一个由肉构成的灵魂，那么祂与我们无关，因为祂没有拯救我们的灵魂。事实上，它也不需要救恩，因为它不是我们的灵魂，既然按照假设，它是一个由肉构成的灵魂。但显然，它已得救。因此，它不是由肉构成的，且它是我们的；因为被拯救的是我们的\Emph{灵魂}，因为它在\Emph{永罚}的危险中。因此我们得出结论：正如在基督内灵魂不是肉，同样祂的肉也不可能由灵魂构成。\par

% structure: p0023 source=h2 target=section reason=relative-heading-rank; base=h2

\section{第十一章 揭露相反的极端。即基督拥有由肉构成的灵魂——有形有体，虽不可见。基督的灵魂，如同我们的，虽穿着肉，却与肉有别。}

但是我们遇到了他们的另一个论点，当我们提出为什么基督在取一个由灵魂组成的肉身时，似乎拥有一个由肉身组成的灵魂？因为天主，他们说，愿意使灵魂对人成为可见的，通过赋予它一个形体的本性，尽管它以前是不可见的；实际上，按其本性，由于这肉身的障碍，它不能看见任何东西，甚至它自己，以致于它是否被生也成了疑问。因此，灵魂（他们接着说）在基督里被造成了有形体的，以便我们能在它经历诞生、死亡以及（更重要的是）复活时看见它。但是，这怎么可能呢？通过肉身，灵魂应该向它自己或向我们展示自己，而在它存在之前，它不可能确定它会提供这种通过肉身展示自己的方式？它接受了黑暗，只是为了能够发光！现在，让我们首先注意这一点：灵魂是否需要以所主张的方式展示自己；其次，\Emph{思考}他们之前的立场是否是灵魂完全不可见的（进一步探究）这种不可见性是其无形体的结果，还是它实际上拥有某种属于自己的身体。然而，尽管他们说灵魂是不可见的，他们却断定它是有形体的，但具有某种不可见的东西。因为如果它没有任何不可见的东西，怎能说它是不可见的呢？但它的存在也是不可能的，除非它有赖以存在的东西。既然它存在，它必然有一个通过它而存在的东西。如果它有这个东西，那一定是它的身体。一切存在的东西都是\OriginalTerm{自成一类}{sui generis}的有形体的存在。只有不存在的东西才缺乏形体。那么，如果灵魂有一个不可见的身体，那愿意使它可见的那位，若使那被认为不可见的部分变得可见，当然会更好地完成祂的工作；因为那样就不会有虚假或软弱在当中了，而这两种缺陷都不适合天主。（但依此假设的情况而言）有\Emph{虚假}，因为祂将灵魂显示为与它实际不同的东西；还有\Emph{软弱}，因为祂无法使它成为它本来的样子。任何一个想要展示一个人的人，都不会用面纱或面具遮住他。然而，这正是对灵魂所做的，如果它被穿上属于别的东西的遮盖，即被转变成肉身。但是，即使灵魂按他们的假设被认为是没有形体的，以至于灵魂，无论它是什么，都应该通过某种\Emph{神秘的}理由力量完全不为人知，只是不是身体，那么在这种情况下，天主的能力并不缺乏——实际上，更符合祂的计划——如果祂在某种新的、与我们所有人都不同的身体中显示灵魂，我们对这种身体会有完全不同的概念（避免了这样的想法）：祂不当地决意要制造一个可见的灵魂而不是一个不可见的灵魂——这无疑是为了他们通过坚持灵魂有人的肉身而引发的那些问题。然而，基督除非作为人，否则不能出现在人中间。因此，请归还基督祂的信德；\Emph{相信}那愿意以人的样式行走在地上的那一位，也显示了一个完全人性的灵魂，不是将它造成肉身，而是以肉身披戴它。\par

% structure: p0025 source=h2 target=section reason=relative-heading-rank; base=h2

\section{第十二章 灵魂的真正功能。基督在祂完美的人性中接受了灵魂，不是为了揭示和解释它，而是为了拯救它。灵魂与肉身一同复活由基督保证。}

好，现在假设灵魂藉肉身显现是基于一个假定——即灵魂显然必须以某种方式显现，因为它对自己和我们来说都是不可知的：这个\OriginalTerm{假说}{hypothesis}中仍然有一个荒谬的区分，它意味着我们自身是与我们的灵魂分离的，而实际上我们的一切都是灵魂。的确，没有灵魂我们就什么都不是；甚至连人的名称都没有，只有尸体的名称。所以，如果我们对灵魂无知，那实际上就是灵魂对自己无知。因此，留给我们审查的仅剩的问题是：灵魂在这件事上是否如此无知于自己，以至它只能以任何可能的方式被认识？在我看来，灵魂是感性的。因此，一切属于灵魂的事物都与感觉有关，一切属于感觉的事物都与灵魂有关。如果我可以强调地使用这个说法，我会说：\Quote{\Emph{Animœ anima sensus est}}——\Quote{Sense is the soul's very soul.}既然赋予一切（有感觉者）感知能力的是灵魂，而且灵魂自己感知所有感觉本身（更不用说它们的性质了），那么它怎么可能没有把感觉作为自己的自然构造而接受呢？如果它不知道自己的特性和它的需要，它又如何能从自然原因的必要性本身知道在特定情况下什么是必要的呢？事实上，认识这一点是每个灵魂的能力；我的意思是，灵魂有一种对自身的实践知识，没有这种自知，任何灵魂都不可能行使自己的功能。我也认为，特别合适的是，人作为唯一的理性动物，应被赋予这样的灵魂，这灵魂本身是卓越理性的，从而使成为理性动物。那么，那使人成为理性动物的灵魂，如果它自己对自己的理性无知，对自己本身无知，它怎么能是理性的呢？然而，它远非无知，它知道自己的创造者、自己的主人和自己的境况。在它学到关于天主的任何事之前，它就称呼天主的名字。在它获得对天主审判的任何知识之前，它就宣称将自己交托给天主。没有比\Quote{死后没有希望}更常听到的话了；然而，根据一个人是善终还是恶终，灵魂会发出怎样的咒诅或祈求啊！这些反思在我们所写的一篇短论\Quote{\WorkTitle{论灵魂的见证}}中有更充分的探讨。此外，如果灵魂从一开始就对自己无知，那么它从基督那里学到的就只能是它自己的性质。它从基督那里学到的不是它自己的形式，而是它的救恩。为此，天主子降下并取了一个灵魂，不是为了让灵魂在基督内发现它自己，而是让基督在灵魂内发现祂自己。因为灵魂的救恩陷入危险，不是由于它对自己无知，而是由于它对天主圣言无知。祂说：\Quote{这生命已显示出来}\BibleRef{若一 1:2}，而不是灵魂。又说：\Quote{我来是为拯救灵魂。}祂没有说\Quote{解释}它。当然，除非灵魂被有形地显示出来，否则我们无法知道灵魂虽然是不可见的本质，却会诞生和死亡。我们确实无知于它将与肉身一同复活。这就是将会发现被基督启示的真理。但即使这一点，祂在自己内启示的方式也不同于某个拉匝禄，他的肉身并非由灵魂组成，他的灵魂也并非由肉身组成。那么，关于灵魂的结构，我们以前无知，现在又得到了什么进一步的知识呢？它有什么属于它的不可见部分需要通过肉身变得可见呢？\par

% structure: p0027 source=h2 target=section reason=relative-heading-rank; base=h2

\section{第十三章。基督的人性。肉身与灵魂两者圆满且无混淆地包含在祂内。}

灵魂成了肉身，好使灵魂成为可见的。那么，肉身是否也同样成了灵魂，好使肉身被彰显呢？如果灵魂是肉身，它就不再是灵魂，而是肉身；如果肉身是灵魂，它就不再是肉身，而是灵魂。那么，在既有肉身又有灵魂之处，它就成了既是这个又是那个。然而，如果它们各自都不是特定的，尽管它们成了既是这个又是那个，那么，至少可以说，这是非常荒谬的：当我们命名肉身时，我们理解为灵魂；当我们指明灵魂时，却解释为我们意指肉身。一切事物都有被误解的风险，即按其自身固有的意义之外的意义被理解，并且，在被那样理解时，失去其本义，如果它们被冠以与它们自然名称不同的称谓的话。名称的忠实保证了属性的正确认识。当这些属性发生变化时，它们就被认为拥有其名称所指示的那些特质。例如，烘烤过的黏土得到了砖的名称。它不再保留标明其先前状态的名称，因为它不再参与那状态。因此，基督的灵魂既然成了肉身，就不能是它已成为的之外的任何东西，也不能再是它从前所是的，因为它确已成了别的。既然我们刚才用了一个比喻，我们将进一步使用它。那么，由黏土制成的我们的水罐，是一个身体，并且有一个名称，当然是指明那个身体；水罐也不能再被称为黏土，因为从前之所是，它已不再是。现在，那不再是（它从前所是）的，也不是不可分离的属性。而灵魂不是不可分离的属性。因此，既然它成了肉身，灵魂就是一个均匀的实体；它也是一个完全单一的存在，不可分割的实体。但在基督身上，我们却发现灵魂和肉身是以简单非比喻性的词汇表达的；也就是说，灵魂被称为灵魂，肉身被称为肉身；没有任何地方灵魂被称为肉身，或肉身被称为灵魂；然而，如果情况果真如此，它们本应被这样（混淆地）命名。\Emph{然而事实却是}，就连\Person{基督}自己也将每个实体分别分开提及，当然，这符合二者属性之间的区别：灵魂单独，肉身单独。祂说：\BibleQuote{\Emph{我的灵魂}甚是忧伤，几乎要死；}以及\BibleQuote{我将要赐给的饼就是\Emph{我的肉身}，（我所要赐的）为世界的生命而赐的。}\BibleRef{若 6:51}现在，如果灵魂本是肉身，那么在基督内只会有一个由肉组成的灵魂，或者一个由灵魂组成的肉身。然而，既然祂保持种类分离，肉身和灵魂，祂显示它们是两个。如果是两个，那么它们就不再是一个；如果不是一个，那么灵魂就不是由肉身组成，肉身也不是由灵魂组成。因为灵魂-肉身，或肉身-灵魂，只是一个；除非祂确实还有另一个灵魂，不同于那曾是肉身的，并且还带着另一个肉身，不同于那曾是灵魂的。但因为祂只有一个肉身和一个灵魂——那\BibleQuote{忧伤至死的灵魂}，和那\Emph{曾是}\BibleQuote{为世界生命而赐的饼}的肉身——两个种类不同的实体的数目保持不变，从而排除了由肉身组成的灵魂这种独特的种类。\par

% structure: p0029 source=h2 target=section reason=relative-heading-rank; base=h2

\section{第十四章 基督没有取天使的本性，而是取了人性。祂来是要拯救人，而不是天使。}

但基督，他们说，取了天使的（本性）。为何如此？此与促使祂成为人的理由相同吗？那么，促使基督行动的动机是那引导祂取得人性的动机。人的救恩是动机，那已丧失之物的恢复。人已丧失；他的恢复成为必要。然而，对于基督取得天使的本性，却不存在这样的原因。因为虽然天使也被判罚在\BibleQuote{为魔鬼及其天使所预备的烈火}中\BibleRef{玛 25:41}，但从未应许给它们恢复。关于天使的救恩，基督从未从父领受任何委托；凡父既未应许也未命令的事，基督不可能承担。那么，祂为何取了天使的本性，若不是（为了让它）作为有力的帮手以执行人的救恩？事实上，天主子独力不能拯救人，而人竟被单独一条蛇推翻了！那么，若有两个筹谋救恩者，且其一需要另一，则不再只有一个天主，只有一个救主。但祂的目的果真要通过一位天使来拯救人吗？那么，为何祂亲自降下，去做那祂要借助天使帮助来加速的事？若借助天使的帮助，为何祂自己也\Emph{来}？若祂打算亲自做一切，为何还要一位天使？诚然，祂被称为\Quote{\OriginalTerm{伟大的谋士天使}{Angel of great counsel}}，即使者，这是表达职务功能而非本性的术语。因为祂必须向世界宣告父的伟大旨意，就是那规定人恢复的旨意。但祂不因此被视为一位天使，如同\Person{加俾额尔}或\Person{弥额尔}。因为葡萄园的主人也差遣祂的儿子到工人那里收取果实，如同差遣仆人一样。然而，儿子不会因此被视为仆人之一，因为祂担任了仆人的职务。那么，若可冒险这样说，我更可以说：儿子实际上是来自父的一位天使（即使者），而并非说在儿子内有一位天使。然而，关于儿子自身已经宣告：\Quote{你使祂稍微逊于天使}；如果祂因成为人，带着血肉和灵魂作为人子而被弄得比天使低微，那么怎会显出祂穿上了天使的本性？然而，正如\BibleQuote{天主的神}和\BibleQuote{至高者的能力}\BibleRef{路 1:35}，祂，那真实的天主和天主子，怎能被视为低于天使呢？但作为承担人性，祂被弄得低于天使；但作为承担天使的本性，祂则同等程度地失去了那低微。这个观点对\Person{厄彼翁}非常合适，他主张耶稣仅仅是一个人，不过是大卫的后裔，而不是天主子；尽管他在某一方面比先知更光荣，因为他宣称在耶稣内有一位天使，如同在\Person{匝加利亚}内一样。只是基督从未说过：\Quote{那在我内说话的天使对我说}\BibleRef{匝 1:14}。此外，基督也从未使用过所有先知常说的那句话：\Quote{上主这样说}。因为祂自己是主，凭自己的权柄公开说话，以\Quote{我实实在在告诉你们}这句话作为前言。还需什么辩论？听依撒意亚用强调的话说：\BibleQuote{不是天使，也不是代表，而是上主亲自拯救了他们}\BibleRef{依 63:9}。\par

% structure: p0031 source=h2 target=section reason=relative-heading-rank; base=h2

\section{第15章 华伦提奴派关于基督肉体具有神性本质的虚构，经圣经检验并驳斥}

\OriginalTerm{瓦伦提努}{Valentinus}确实，凭借其异端体系，可以自圆其说地为基督设想一种属神的肉体。任何拒绝相信那肉体是人的肉体的人，都可以假装它是什么都行，因为（而且这句话适用于所有的\Emph{异端者}），如果它不是人的，也不是由人所生，我就看不出基督本人称自己为人和人子时，说的是什么实体，祂\Emph{说道}：\BibleQuote{如今你们却想杀我，一个告诉你们真理的人；}\BibleRef{若 8:40}和\BibleQuote{人子也是安息日的主。}\BibleRef{玛 12:8}因为依撒意亚正是论及祂写道：\BibleQuote{受苦的人，熟悉如何承受软弱；}耶肋米亚：\BibleQuote{祂是一个人，谁能认识祂？}达尼尔：\BibleQuote{乘着云彩而来的人子。}\BibleRef{达 7:13}宗徒保禄同样说：\BibleQuote{那人基督耶稣是人与天主之间的唯一中保。}\BibleRef{弟前 2:5}伯多禄在\WorkTitle{宗徒大事录}中也论及祂确是真有人性（他说）：\BibleQuote{耶稣基督是天主在你们中间认可的人。}\BibleRef{宗 2:22}这些段落本身应当足以作为定论性的证据，证明基督拥有来自人的肉体，而非属神的；并且祂的肉体不是由灵魂构成的，也不是由星辰物质构成的，也不是虚幻的肉体；（毫无疑问，它们本是足够的），只要异端者们能够摆脱他们所有好争辩的热情和诡计。因为我曾在瓦伦提努可悲派系的某位作者的作品中读到，他们从一开始就拒绝相信人性及地上的实体是为基督创造的，免得主被认为低于那些不是由地上的肉体形成的天使；由此，如果祂的肉体像我们的一样，那么它也必然不是由圣神、也不是由天主所生，而是由人的意志所生。再者，为什么它应该是由不朽的（种子）所生，而不是由可朽的？又，既然祂的肉体已经复活并升天，为什么我们的肉体（与祂相似）不是也立刻被提升？或者，为什么祂的肉体（既然与我们相似）不照样归于尘土并腐朽？这样的反对意见甚至连外邦人也经常互相争论。天主之子竟沦落到如此卑贱的地步吗？再者，如果祂复活作为我们望德的先例，为什么我们自身却不认为发生类似的事是可取的？这样的观点对于外邦人并非不当，对于异端者也是合宜和自然的。因为，他们之间究竟有什么区别，除了外邦人不相信，却相信；而异端者相信，却不相信？然后，他们又读到：\BibleQuote{祢使祂稍微逊于天使；}因此他们否认那自称\BibleQuote{不是一个人，而是一条虫}的基督的卑微本性；祂也没有\BibleQuote{俊美形貌，但祂的容貌卑贱，被人轻视，多受痛苦，熟悉如何承受软弱。}在这里，他们发现一个人与一位神性者相混合，因此他们否认其人性的身份。他们相信祂死了，却主张一个已经死了的存在是由不朽的实体所生；仿佛可朽与死亡是两回事！但是我们的肉体也应该立即复活。等待片刻。基督尚未制服祂的敌人，以便能与祂的朋友们一起战胜它们。\par

% structure: p0033 source=h2 target=section reason=relative-heading-rank; base=h2

\section{第十六章 基督的肉体在本质上与我们相同，只是无罪。\OriginalTerm{罪的肉体}{Carnem Peccati}与\OriginalTerm{肉体的罪}{Peccatum Carnis}的区别：基督所废除的是后者。第一个亚当的肉体，与第二个亚当的肉体一样，并非由人的种子所受，虽然如同我们自己的肉体一样完全是人的，后者源于它。}

那位著名的亚历山大，同样受好辩之心所煽动，带着异端性情的真实风范，竟公然反对我们；他声称我们主张基督穿上源自尘世的血肉，为的是祂能亲在自己身上消除罪恶之肉。然而，即使我们确持此见，我们也能够加以辩护，完全避开他所归咎于我们的狂妄蠢行——他以为我们竟认为基督自身的那同一血肉被当作有罪而消亡了；因为我们在公开场合宣认我们的信仰，即那血肉正坐于圣父在天之右，并且我们还宣告它将从那里带着圣父全部荣耀的威仪再次来临：因此，我们既不可能说它已经消亡，正如我们不可能坚持它有罪而归于虚无，因为它之中未曾有过任何过错。此外，我们坚持认为，在基督里被消除的并非\OriginalTerm{罪恶之肉}{carnem peccati}，\Quote{罪恶之肉}，而是\OriginalTerm{肉中之罪}{peccatum carnis}，\Quote{肉中之罪}——不是那质料，而是其状况；不是那实体，而是其瑕疵；并且我们以宗徒的权威断言此事，他说：\Quote{祂在肉中消灭了罪。}在另一处经文里，他又说基督\Quote{在罪恶之肉的相似中}，然而，并非祂取了一个\Quote{血肉的相似}，仿佛是一个身体的虚影而非实体；而是他要我们理解为相似于那犯过罪的血肉，因为基督的血肉本身并未犯过罪，却相似于那犯过罪的血肉——在本性上相似，但在亚当所承受的败坏上并不相似；因此我们也肯定，在基督里存在和人在其中本性有罪的血肉相同的血肉。所以，在血肉中，\Emph{我们说}罪恶已被消除，因为在基督里那同一血肉保持了无罪，而在人身上它却未保持无罪。现在，倘若基督不是在具有罪之本性的那个血肉中消除罪恶，那就不符合祂在血肉中消除罪恶的目的，也无助于祂的荣耀。因为倘若祂在某种更好的、具有另一种甚至无罪本性的血肉中除去了罪的污点，那当然算不上什么奇事！那么，你便说：如果祂取了我们的血肉，基督的血肉就是有罪的。然而，不要用一个晦涩难解的概念来束缚这十分明了的含义。因为祂在穿上我们的血肉时，使之成为自己的；在使之成为自己的同时，便使之成为无罪的。但必须向所有因基督的血肉并非来自人间父亲的种子而拒绝相信我们血肉在基督里的人进一言：让他们记住，亚当自己也是在没有人间父亲种子的情况下领受了我们这血肉。正如大地未经人间父亲种子而变化为我们这血肉，同样，天主圣子也完全可以在没有人间父亲作为工具的情况下，将那同一血肉的本体取归己有。\par

% structure: p0035 source=h2 target=section reason=relative-heading-rank; base=h2

\section{第17章　第一位与第二位亚当在其血肉来源方面的处境相似。在厄娃与童贞玛利亚之间亦愉快地描绘出一个类比。}

但是，我们暂且撇开亚历山大及其三段论——他在讨论中如此乖僻地运用它们——以及瓦伦提努斯的颂歌，他竟以绝顶自信将其掺入，作为某位可敬作者的著作；让我们把探究局限于一点：基督是否从童贞女取了肉躯？这样，如果祂从母胎取得了其质料，我们就能得出确证，证明祂的肉躯是人的；尽管我们立刻就有明显的证据，表明祂肉躯的人性特征，从作为人的名号和描述，从其构造的本性，从其感官系统，以及从其遭受死亡。现在，首先有必要说明，圣子由童贞女所生先前有何理由。那将要祝圣一种新生育秩序者，祂自己必须以一种新颖的方式诞生，关于此事，依撒意亚曾预言主将亲自赐予征兆。那么，这征兆是什么？\BibleQuote{看，一位童贞女将怀孕生子。}\BibleRef{依 7:14}于是，一位童贞女确实怀孕生了\BibleQuote{厄玛奴耳，天主与我们同在。}\BibleRef{玛 1:23}这是新的诞生：一个人在天主内诞生。而在这人内，天主诞生，取了古族的肉躯，但未借助古种，为要以新种，即以属灵的方式，改造它，并除去其一切古污以洁净它。但整个这新生，如同在其他一切情形中一样，已在古型中预表：主藉一种安排而作为人诞生，其中童贞女为媒介。大地尚处于童贞状态，未经人手耕种，尚无种子撒入犁沟，那时，如我们被告知的，天主用地上的土造了人，成为有灵的生命。\BibleRef{创 2:7}既然第一位亚当这样被介绍给我们，那么合理的推论是，第二位亚当同样，如宗徒告诉我们的，是由天主从土地中形成，成为赋予生命的灵——换句话说，出自尚未被任何人所玷污的肉躯。但为了不失去任何从亚当之名支持论证的机会，基督为什么被宗徒称为亚当呢？除非是因为，作为人，祂出自那尘世的根源。而连理性在此也维持同样的结论，因为正是藉着相反的行动，天主恢复了被魔鬼夺去的自己的肖像和模样。因为当厄娃还是童贞女时，那诱骗的话语已潜入她的耳朵，这话语将建造死亡的殿宇。同样，必须将天主的圣言引入一位童贞女的灵魂，这圣言将搭建生命的殿宇；这样，由这性别所导致毁灭的，可由同一性别恢复而得救恩。犹如厄娃相信了蛇，玛利亚也相信了天使。一个因相信而造成的过犯，另一个因相信而抹除。但（或许会说）厄娃并未因魔鬼的话而在腹中怀孕。唉，她无论如何怀了孕；因为魔鬼的话后来成了她的种子，使她像被逐者一样怀孕，并在痛苦中生产。她确实生了一个杀弟的魔鬼；而玛利亚相反，生了那一位，祂将在某一日为以色列确保救恩，祂自己按肉身的兄弟，也是杀害祂自己的人。因此，天主将祂的圣言降下到童贞女的胎中，作为那位善兄弟，他要抹去恶兄弟的记忆。因此，基督必须带着那种\Emph{肉躯}的状态出来，以拯救人类，人类自从被定罪以来就进入了那种状态。\par

% structure: p0037 source=h2 target=section reason=relative-heading-rank; base=h2

\section{第十八章 天主圣三第二位格摄取我们完美人性的奥迹。祂在这里，如同在其他许多地方一样，被称为\OriginalTerm{圣神}{Spirit}。}

为要给予一个更简单的回答：\Person{天主子}不宜由人的父亲之种子所生，免得祂若完全是人的儿子，便不再是天主子，并比\Quote{一个\Person{撒罗满}}或\Quote{一个\Person{约纳}}更无长处。\BibleRef{玛 12:41} 如\Person{厄俾翁}以为我们应如此相信祂。因此，为使那已是天主子——即由天主圣父的种子，也就是说，圣神——也为人子，祂只需要取得血肉，取人的血肉而不需人的种子；因为对于那已持有天主种子者，人的种子是不必要的。所以，正如祂在由童贞女出生之前，能无人的母亲而拥有天主为父，同样，在祂由童贞女出生之后，也能无人的父亲而拥有女子为母。总之，祂与天主同为人，因为祂以天主的圣神而拥有人的血肉——我所说的\Emph{血肉}，非由人的种子而来；\Emph{圣神}，则由天主的种子而来。既然天主的救恩计划关乎祂圣子的旨意，要求祂应由童贞女所生，为何祂不能从童贞女那里接受祂所携带的身体——那身体是来自童贞女的？因为（他们说）祂从天主那里领受了别的什么，因为\BibleQuote{圣言成了血肉}。\BibleRef{若 1:14} 这一陈述清楚地显明了那成为血肉的是什么；不可能是别的，只能是圣言成了血肉。那么，圣言是由血肉成了血肉，还是由（神性的）种子本身成了血肉，圣经必须告诉我们。然而，圣经对于除了那成为血肉者以外的一切都缄默不语，且只字不提它是由什么而成的，因此必须认为它暗示圣言是由别的东西、而非由它自己成了血肉。如果不是由它自己，而是由别的东西，我们还能更合宜地假定圣言成为血肉的是从哪一样东西，而不是从那血肉——在其中祂屈从于那救恩计划——呢？（我们还有一个相同的结论的证明）就是主自己郑重而清楚地宣称：\BibleQuote{那由血肉生的就是血肉}，\BibleRef{若 3:6} 因为它是由血肉生的。但若祂在这里只是说普通人，而非祂自己（正如你们所主张的），那么你们必须完全否认基督是人，并坚持说\Emph{人性}并不适合祂。然后祂接着说：\BibleQuote{那由圣神生的就是神}，\BibleRef{若 3:6} 因为天主是神，而祂是由天主生的。这描述其实对祂比对相信祂的人更适用。但如果这经文确实指祂，为何前一句不也指祂呢？因为你们不能将两者的关系分割，并将这一句归于祂，而前一句归于所有其他人，尤其是你们并不否认基督兼具血肉和圣神两性。此外，既然祂既具血肉又具圣神，祂在论及祂自己所携带的两性的状况时，不可能被理解为断定圣神确属祂自己，而血肉却不属祂自己。所以，既然祂由圣神而来，祂就是天主圣神，并由天主而生；同样，祂也由人的血肉而生，作为人在血肉中受生。\par

% structure: p0039 source=h2 target=section reason=relative-heading-rank; base=h2

\section{第十九章　基督按其神性，作为天主的圣言，成为血肉，不是藉属肉体的受孕，也不是藉血肉和人的意愿，而是藉天主的旨意。基督的神性自甘下降于童贞女的子宫。}

那么，这段经文是什么意思：\Quote{不是由血、也不是由肉情的意愿、也不是由人的意愿，而是由天主而生的？}\BibleRef{若 1:13} 在驳斥那些窜改经文的人之后，我将更充分地引用这段经文。他们声称，原文是这样写的（复数形式）：\Quote{\Emph{他们}不是由血、也不是由肉情的意愿、也不是由人的意愿，而是由天主而生的，}好像是指那些前文提到的\Quote{信祂之名的人，}以便指出他们私自占有的那奥秘的、属灵的被选之种的存在。但这怎么可能呢？所有信主之名的人，由于人类共同的本原，都是由血、肉情的意愿和人而生的，正如华伦提努本人一样。经文用的是单数，指主而言：\Quote{祂是由天主而生的。}这十分恰当，因为基督是天主之圣言，与圣言同在的是圣神，藉着圣神有天主的能力，以及其他一切属于天主的。然而，作为血肉，祂不是由血、也非肉情的意愿、也非由人，因为是天主愿意圣言成了血肉。这否认了那对我们所有人而言自然的诞生，的确归于血肉，而非归于圣言，因为祂是作为血肉如此诞生的，而不是作为圣言。然而，当经文实际上否认祂是由肉情的意愿而生时，怎么没有也否认（祂是由）血肉的本体（而生）呢？因为当它否认祂是\Quote{由血而生}时，并没有否认血肉的本体，而只是否认了种子的物质，众所周知，那是通过蒸腾转化为女人血液之\OriginalTerm{凝块}{coagulum}的热血。在奶酪中，正是通过凝结，乳状物质获得了那种稠度，通过注入凝乳酶而凝固。因此我们明白，所否认的是主藉由性交而诞生（正如\Quote{人的意愿和肉情的意愿}这话所暗示的），而非祂从女人子宫中的诞生。再者，为什么要如此强调祂不是由血、不是由肉情的意愿、也不是由（意愿）人生，如果不是因为祂的血肉是那样的，以致没有人会对它是从性交而生这一点有任何怀疑呢？同样，虽然否认祂从这样的结合而生，但经文并没有否认祂是由\Emph{真实}血肉所生；相反，它肯定了这一点，因为它没有否认祂在血肉中的诞生，正如它否认祂从性交而生一样。请问我，如果圣神不是为从子宫取得血肉，那祂为何要降入妇人的子宫呢？因为祂本可以不通过这一过程就成为\Emph{属灵}的血肉——事实上，没有子宫比在子宫中更简单。如果祂不从其中带出什么，祂就没有理由将自己封闭在其中。然而，祂降入子宫并非没有理由。因此，祂从那里取得了（血肉）；否则，如果祂从那里什么也没有取得，祂降入其中就是毫无理由的，尤其如果祂打算成为那种并非从子宫而来的血肉，即属灵的血肉。\par

% structure: p0041 source=h2 target=section reason=relative-heading-rank; base=h2

\section{第20章 基督由童贞女诞生，出自她的实体。根据某些圣经经文所提示的祂由一位人类母亲真实而精确诞生的生理事实。}

然而，你竟采用何等狡计，企图剥夺介词\Emph{ex}（\Emph{从}）作为介词的固有力量，并用另一个不曾在整部圣经中出现过的含义来替代它！你说祂是\Emph{通过}一位童贞女而生，而非\Emph{从}一位童贞女而生；是\Emph{在}子宫中，而非\Emph{从}子宫中出生，因为天使在梦中向若瑟说：\Quote{那在她内（不是从她）所生的，是出于圣神。}\BibleRef{玛 1:20} 但事实上，若他本意是\Quote{从她}，他必定会说\Quote{在她内}；因为凡是从她而出的，也就在她内。因此，天使所说的\Quote{在她内}一词，与\Quote{从她}这一短语具有完全相同的含义。不过，幸运的是，玛窦在追溯主从亚巴郎到玛利亚的世系时也说：\Quote{雅各伯生若瑟，玛利亚的丈夫，基督是由她（\OriginalTerm{从她}{ex qua}）所生的。}\BibleRef{玛 1:16} 此外，保禄也堵住了这些批评者的口，他说：\Quote{天主派遣了自己的圣子，由女人所生。}\BibleRef{迦 4:4} 他指的是\Emph{通过}女人，还是\Emph{在}女人内？不仅如此，为了更加强调，他用了\Quote{\Emph{所造}}这个词，而不是\Emph{所生}，尽管用后者更为简单。但藉着说\Quote{所造}，他不但证实了\Quote{圣言成了血肉}\BibleRef{若 1:14} 这句话，也肯定了那由童贞女所造的血肉的实在性。我们在\WorkTitle{圣咏}上也将得到支持，当然不是背教者、异端者和柏拉图主义者瓦伦提努的\WorkTitle{圣咏}，而是达味的\WorkTitle{圣咏}，这位最杰出的圣人和众所周知的先知。他向基督歌唱，并且透过他的声音，基督也实在歌唱了祂自己。那么，请听基督主对天主圣父所说的话：\Quote{是祢从我母胎中牵引了我。}这是第一点。\Quote{祢是我从我母亲乳房的希望；我从母胎就仰赖了祢。}这是第二点。\Quote{祢是我从我母亲腹中的天主。}这是第三点。\newline{}现在让我们仔细留意这些段落的意思。\Quote{祢牵引了我，}祂说，\Quote{从母胎中。}那么，什么是被\Emph{牵引}的呢？岂不就是那附着之物、那被牢牢系在某个从之被牵引以便分离的东西上之物？如果祂不附着于母胎，怎能被从母胎中牵引出来？若那附着于母胎之物被从其中牵引出来，祂又如何能附着于它呢？除非祂在子宫中的整个期间，都藉着脐带与其相连，作为祂的起源，而脐带则从子宫把生长滋养传递给祂？即使一种异质与另一种异质混合，它也会如此完全地与它所混合之物结合，以致当它从其中被分离出来时，会带走它所撕裂的物体的一部分，仿佛是由于各部分之前相互结合与生长所产生的撕裂。但是祂所提到的祂\Quote{母亲的乳房}是什么呢？无疑就是祂所吸吮的乳房。助产士、医生和博物学家能根据女性乳房的特性告诉我们，乳房通常除了在子宫受孕时之外，是否会在其他时间流乳？此时，静脉将下部的血液输送到\OriginalTerm{乳房}{mamilla}，并在输送过程中将分泌物转化为营养丰富的乳汁。由此导致哺乳期间经期停止。但是，如果圣言是凭己意成为血肉，没有与子宫有任何接触，没有母体的子宫以通常的功能和滋养对其产生作用，那么乳泉怎能（从子宫）传送到乳房呢？因为（子宫）只有实际拥有\Emph{适当的实体}才能进行这种转化。但它不可能有血液转化为乳汁，除非它也拥有血液的成因，也就是说，它的血肉（通过出生）与\Emph{母亲的子宫}分离。\newline{}现在很容易看出基督由童贞女所生的新颖之处何在。这仅仅在于，祂是（被生）出于一位童贞女，并以我们所指出的真实方式，好使我们的重生能拥有童贞的纯洁——藉着基督在灵性上洁净一切污秽，而基督自己本身就是一位童贞者，甚至在其血肉中也是如此，因为祂是\Emph{从}童贞女的血肉所\Emph{生}的。\par

% structure: p0043 source=h2 target=section reason=relative-heading-rank; base=h2

\section{第二十一章 天主的圣言除非在童贞女的子宫中并出于她的实体，否则没有成为血肉。祂通过祂的母亲是出自她伟大祖先达味的后裔。在旧约和新约中祂都被描述为\Quote{达味腰间的果实}。}

既然他们争辩说，这（基督诞生的）新意在于：圣言降生成人时没有人的父亲的种子，因而也没有童贞母亲的肉身参与此事，那么，为什么这新意不更应当局限于：祂的肉身虽非由种子而生，却仍出自肉身？我愿更深入地探讨这个问题。\BibleQuote{看，}他说，\BibleQuote{童贞女将怀孕在母腹中。}\Emph{怀}什么？我问。当然是圣言，而不是人的种子，而且无疑是为了生一个儿子。\BibleQuote{因为，}他说，\BibleQuote{她要生一个儿子。}因此，既然受孕的行为是她自己的，那么她所生的也是她自己的，尽管受孕的原因不是她的。另一方面，如果圣言独自成为肉身，那么祂就自己受孕并自己生产，那么预言就成荒谬了。因为在这种情况下，童贞女既\Emph{没有}受孕，也\Emph{没有}生产；因为她从圣言受孕所生的一切，都不是她自己的肉身。但这是唯一落空的预言吗？难道天使的宣告不也被颠覆了吗，即童贞女要\BibleQuote{在母腹中怀孕并生一个儿子？}\BibleRef{路 1:31}事实上，一切宣告基督有母亲的圣经岂不也都落空？因为除非祂曾在她腹中，她怎能成为祂的母亲？但祂从她腹中什么也没有得到，那腹中本可使她成为母亲。这样的名号，一个陌生的肉身不应当\Emph{承担}。没有肉身能称自己出自母亲的腹中，除非它本身就是那腹中的产物；也没有任何东西能是那腹中的产物，如果它的诞生只凭自己。因此，连依撒伯尔也必须沉默，尽管她怀中的是先知性的婴孩，这婴孩已意识自己的主，并且充满了圣神。\BibleRef{路 1:41}因为她说的没有道理，\BibleQuote{我主的母亲到我这里来，这事怎会临到我？}如果玛利亚怀耶稣不是作为她的儿子，而只是作为一个陌生人，她又怎能说：\BibleQuote{你腹中的果实是有福的？}这腹中的果实是什么？它没有从腹中接受种子，没有在腹中扎根，不属于那腹的主人，而这无疑正是腹中真实的果实——甚至基督？如今，既然祂是从叶瑟的根发出的枝条的嫩芽；再者，叶瑟的根是大卫的家族，根的枝条是玛利亚，她出自大卫，枝条的嫩芽是玛利亚的儿子，称为耶稣基督，祂不也是果实吗？因为嫩芽就是果实，因为通过嫩芽并从嫩芽，每个产物从初始状态进展到完美的果实。那么怎样？\Emph{他们}否认果实有嫩芽，嫩芽有枝条，枝条有根；因此，根通过枝条未能为自己获得那来自枝条的特殊产物，即嫩芽和果实；因为事实上，家谱中的每一步都是从最近追溯至最先，所以如今一个众所周知的事实是，基督的肉身不仅与玛利亚不可分离，而且通过玛利亚与大卫，通过大卫与叶瑟也不可分离。因此，\Quote{这果实，}因此，\Quote{出自大卫的腰，}即出自他肉身的后裔，天主向他起誓，\Quote{他应兴起坐在他的宝座上。}如果\Quote{出自大卫的腰，}那么祂更出自玛利亚的腰，因为借着她，祂在\Quote{大卫的腰中。}\par

% structure: p0045 source=h2 target=section reason=relative-heading-rank; base=h2

\section{第二十二章 新约圣经，甚至在其第一节经文，就为基督真实的肉身作证。凭借这肉身，祂被纳入大卫、亚巴郎和亚当的人类谱系。}

他们或许可以抹煞魔鬼的见证，这些魔鬼曾宣称耶稣是达味之子；但不论这见证如何不足取，宗徒们的见证却是永远无法抹去的。首先，玛窦，这位最忠实的福音编年史家，因为他是主的同伴，他为了清楚地向我们指明基督的血缘出身，就以这样的方式开始\Emph{他的福音}：\BibleQuote{耶稣基督、达味之子、亚巴郎之子的族谱。} \BibleRef{玛 1:1} 出自如此本源的性体，并且按顺序逐渐下降到基督的诞生，这里所描述的，岂不就是亚巴郎和达味的肉身，一步一步地传递到童贞女，最终引进了基督，——不，由童贞女产生了基督自己吗？然后，又有保禄，他同时是同一福音的门徒、导师和见证人；作为同一基督的宗徒，他也证实基督\Quote{按肉身是达味的后裔}——因此这也是他自己的。所以，基督的肉身是达味的后裔。既然祂是因玛利亚的肉身而成为达味的后裔，那么祂也是因达味的后裔而成了玛利亚的肉身。无论你如何曲解这句话，祂要么是因达味的后裔而成了玛利亚的肉身，要么是因玛利亚的肉身而成了达味的后裔。同一个宗徒在宣告基督是\Quote{亚巴郎的后裔}时结束了整个讨论。如果是亚巴郎的后裔，那么当然更是达味的后裔，因为达味是更近的\Emph{祖先}。因为，当保禄在亚巴郎身上展开对万民的应许祝福时，\Quote{地上万民都要因你的后裔蒙福}，他接着又说：\Quote{他不是说\InnerQuote{后裔们}，如同指许多人；而是指一个人\InnerQuote{你的后裔}，就是基督。} \BibleRef{迦 3:8} 当我们阅读并相信这些时，我们应当且能够承认基督拥有怎样的肉身呢？当然只能是亚巴郎的肉身，因为基督是\Quote{亚巴郎的后裔}；只能是叶瑟的肉身，因为基督是\Quote{叶瑟的树干}的苗芽；只能是达味的肉身，因为基督是\Quote{达味腰间的果子}；只能是玛利亚的肉身，因为基督出自玛利亚的胎；而更高一层，只能是亚当的肉身，因为基督是\Quote{第二亚当}。因此，结论是：他们要么坚持那些祖先拥有属神的肉身，以便基督能承袭同样的实体状况，要么承认基督的肉身不是属神的，因为它并非来自属神血统的起源。\par

% structure: p0047 source=h2 target=section reason=relative-heading-rank; base=h2

\section{第23章。西默盎的\Quote{该受反对的记号}，应用于异端者对基督真正诞生的反驳。异端者的一个悖论转为支持公教真理。}

我们承认，西默盎的先知预言实现了，这预言是他论到刚出生的救主所说的：\BibleQuote{看，这个\Emph{孩子}已被立定，为使以色列中许多人跌倒和兴起，并成为反对的记号。} \BibleRef{路 2:34} 这记号（此处所指的）是基督诞生的记号，按照依撒意亚：\BibleQuote{因此主要亲自给你们一个\Emph{记号}：看，童贞女要怀孕生子。} \BibleRef{依 7:14} 我们于是发现那该受反对的记号是什么——童贞玛利亚的受孕和分娩，关于这个，这些诡辩者说：\Quote{她是童贞女，又不是童贞女，生了，又没有生；} 仿佛这种语言，如果确实必须说出来，连我们自己使用不是更合适吗！因为\Quote{她生了}，因为她产生了自己肉身的后代，而\Quote{她没有生}，因为她不是从丈夫的种子产生了祂；她是\Quote{童贞女}，就远离丈夫而言，而\Quote{不是童贞女}，就她生孩子而言。\Emph{然而，并没有异端者所假装的那种推理等同：换言之}不能因为\Quote{她没有生}，就推论出那\Quote{不是童贞女}的\Quote{仍是童贞女}，即使她成为母亲却没有自己胎中的果实。但对我们来说，没有暧昧，没有扭曲的双重意义。光是光，暗是暗；是就是是，非就是非；\Quote{多余的，便是出于邪恶。} \BibleRef{玛 5:37} 那生了（确实）生了；尽管她受孕时是童贞女，她生产\Emph{她的儿子}时却是妻子。现在，作为妻子，她处在\Quote{开胎}的法律之下，其中生产男性是凭借丈夫的合作与否并不重要；是同一个性别开了她的胎。的确，她的胎正是因此，关于其他人也写着：\Quote{凡开胎的男性，都该称为归于主的圣者。} 因为除了圣子，谁是真正圣善的呢？除了祂开了\Emph{一个关闭的胎}，谁真正开了胎呢？但在所有情况中，是婚姻开了胎。因此，\Emph{童贞女的胎}是特别被开的，因为它特别关闭着。的确，她更应该被称为不是童贞女，而不是童贞女，仿佛一跃成为母亲，在她成为妻子之前。关于这一点还需要说什么呢？因为宗徒正是在这个意义上宣称，圣子不是生于童贞女，而是\Quote{生于女人}；他在那声明中承认了婚姻之后\Quote{开胎}的状况。我们在厄则克耳中读到\Quote{一头小母牛生了，又没有生。} 现在，你要看看，难道圣神不是预见到你们将来关于玛利亚的胎的争论，而在这段文字中给你们留下了记号吗？否则祂不会，违背祂在这位先知中通常的简朴文风，说出这样一句有疑义的话，\Emph{尤其}当依撒意亚说：\BibleQuote{她要怀孕生子。} \BibleRef{依 7:14}\par

% structure: p0049 source=h2 target=section reason=relative-heading-rank; base=h2

\section{第二十四章 在先知圣经的各个段落中对各种异端者的神圣谴责。那些攻击独一主耶稣基督（既是天主又是人）之真道的人，如此被定罪。}

因为当依撒意亚针对我们的异端者发出谴责，尤其是在他的\Quote{祸哉，那些称恶为善、以暗为光的人}\BibleRef{依 5:20}中，他当然是在标记你们中间那些在所用言辞中不保持其真实意义的亮光的人（注意）——即\Emph{灵魂}应当只指所谓的那一个，\Emph{肉身}仅仅指向我们明认的那一个，\Emph{天主}无非是所宣讲的那一位。既然以先知的目光有了马尔强，他说：\Quote{我是天主，再没有别的；除我以外没有天主。}\BibleRef{依 45:5}当他在另一处同样说：\Quote{在我以前没有神}\BibleRef{依 46:9}，他击打了瓦伦提尼安永世那不可测的宗谱。同样，圣经中对厄俾翁的答复是：\Quote{不是由血气，也不是由肉情，也不是由人意，而是由天主生的。}同样，在经文\Quote{即使从天上来的天使，若给你们宣讲我们所宣讲的福音以外的其他福音，当受诅咒}\BibleRef{迦 1:8}中，他提醒注意阿佩莱斯的童贞女友斐卢墨涅的狡猾影响。\BibleQuote{那否认基督降生成人肉身的，确实是假基督。}\BibleRef{若壹 4:3}通过宣告祂的肉身是纯粹而绝对真实的，按其自身性质的本义接受，圣经对所有在其中作区分的人予以打击。同样，当它定义基督本身只有一个时，它动摇了那些展示多种基督、使基督为一位而耶稣为另一位的人的空想——他们描述一位从人群中逃出，另一位被他们拘留；一位在孤山上向三个同伴显现，在云彩中荣耀覆盖，另一位为普通人与众人交往；一位宽宏大量，另一位胆怯；最后，一位忍受死亡，另一位复活，通过这个事件他们也主张自己的复活，不过是在另一种肉身中。幸好，那忍受痛苦的\Quote{将从天上再来}\BibleRef{宗 1:11}，并且所有人都将看见那从死者中复活的一位。那些钉死祂的人也将看见并承认祂；也就是说，祂的肉身本身，就是他们对其施暴的那个肉身，没有它祂自己既不可能存在也不可能被看见；因此那些断言祂的肉身如空鞘般无感觉地坐在天上、基督已从中抽身的人，以及那些主张祂的肉身和灵魂完全相同、或者只有灵魂存在而肉身不再活着的人，必将羞愧脸红。\par

% structure: p0051 source=h2 target=section reason=relative-heading-rank; base=h2

\section{第二十五章 结论。本篇论文构成另一部著作\WorkTitle{论肉身的复活}的序言，证明那真实诞生、死亡并复活的肉身的真实性。}

但这在当前主题上就足够了；因为我想此时已经充分证明了基督的肉身既由童贞女所生，且其本质是人的。单凭这一讨论本已足够，无需碰触从不同角落提出的孤立意见。然而，我们已将这些意见引向考验，既基于支持它们的论据，也基于所引用的圣经，而且我们是\OriginalTerm{充足地}{ex abundanti}这样做的；因此，通过展示基督的肉身是什么、来自何处，我们也针对所有反对者预先决定了那肉身不是什么。然而，我们自己的肉身的复活将不得不在另一篇小论文中加以维护，并以此结束当前这篇作为一般序言的作品，它将为即将讨论的主题铺平道路，既然如今已清楚在基督里复活的那身体是哪一种身体。\par

\SourceRecord{Translated by Peter Holmes.；Ante-Nicene Fathers；Vol. 3.；Edited by Alexander Roberts, James Donaldson, and A. Cleveland Coxe.；Buffalo, NY: Christian Literature Publishing Co.,；1885.；<http://www.newadvent.org/fathers/0315.htm>.}
